All data are simulated. The study has one dynamical system (Lorenz-63), one sampling interval ($\Delta t=0.01$), one trajectory length (10 time units), one library (degree-2 polynomials, which contains the true model) and 20 trajectories per cell, so rate differences up to about 0.25 are not reliably resolved (Sec.~\ref{sec:results}) and the $p$-values are uncorrected. All five methods are reimplemented in numpy with small hyperparameter grids whose edge is selected in most cells, so a wider grid might change the ranking. Tuning is oracle-assisted, which inflates absolute rates and is not available in practice. Five tuning trajectories do not determine the threshold: ties in support rate and median error decided it in 19 of 30 cells (in two of them, FD at 5\% and spline at 10\%, between thresholds that give different models), the rule that breaks them was changed after the first test results were seen, and the ordering of the systems below 5\% noise follows that rule. Noise is i.i.d. Gaussian with a level known to the tuner. The weak form uses one test-function family and window count; other weak-SINDy variants, ensembles and implicit forms were not compared. A full study would add systems, sampling rates, noise models, tuning without ground truth on a larger tuning set, and more trials per cell.
